% Ad M. Brutum 1.10 --- headnote
% ad-brutum-01-10, mid-July 43 BC, Rome

\textit{Cicero to M. Brutus, from Rome, mid-July 43 BC
--- Perseus dateline \textit{Scr. Romae med. m. Quint.
a. 711 (43)}. One of the longest and most clear-sighted
of Cicero's late letters: a survey of the political
collapse already in motion. Antony has been defeated at
Mutina in April; both consuls (Hirtius, Pansa) are dead;
the Senate's victory has slipped (``so many faults of
Brutus's followed that in a way the victory slipped out
of our hands''); Lepidus has joined the remnant of
Antony's forces and reopened the war; the young Caesar
(Octavian), guided so far by Cicero, has been worked on
by Caesarian agents and now publicly intends to claim
the consulship. Section 3 is the most precise
contemporary statement of where the danger is coming
from: ``power is now lodged in violence and arms,'' and
``each man demands for himself as much power in the
commonwealth as he has strength.''}

\textit{The structural argument of the letter is the
brief Cicero has been making since spring. He himself,
having tried to escape Italy in the months after
Caesar's assassination, was called back by Brutus's
edicts and by Brutus's exhortation in person at Velia.
He shook Antony with no guard at his side. He
strengthened the guard Octavian's force offered. Now
the guard is at risk of failing: ``if the counsels of
the impious shall prove stronger than mine, or if the
weakness of his age cannot bear the weight of these
affairs, every hope is in you.'' The closing lines of
section 5 are the personal note ordinarily kept off the
page: the standing-ground of hope is on Brutus's parade-
ground, not in Rome; and if it should turn out
otherwise, ``I shall grieve on behalf of the
commonwealth, which ought to have been immortal; for
myself --- how little is left to me?'' Within four
months the answer to that question would be no time at
all.}
